\cotefile{162-5}
\scanpage{46}
$g$ manifestement se prolonge en $\overline{g} : \overline{T} \to X\mathrm{x}Y$ . Cela nous ramène à prouver
que \inserted{le composé} $\unsure{F'}\overline{T}' \xrightarrow{\alpha_i} \unsure{F'}Z'_i \xrightarrow{\alpha_{i*}} \unsure{F'}\overline{T}'$ ne dépend pas de $i=0,1$. Ecrivons $X$ au lieu de $\overline{T}'$,et
considérons le diagramme cartésien d'inclusions
\diagram{\begin{tikzcd}[ampersand replacement=\&]
X_s \arrow[r, "\alpha_s"] \arrow[d, "\alpha_s"'] \& X \arrow[d, "{\Gamma \unsure{\overset{\Delta}{=}|=} (p,\mathrm{id}_X)}"] \\
X \arrow[r, "i_s"] \& S^1\mathrm{x}X
\end{tikzcd}}
,

on a alors $\alpha_s^*\alpha_{s*}=\Gamma^*i_{s*}$ , qui ne dépend pas de $s$ puisque $i_s$ n'en
dépend pas à homotopie près. - Le reste est sans doute l'AQT.

Sur la strcture des $B.(X)$. Ce sont tous des algèbres graduées
sur $B.(\mathrm{pt})=\Lambda$ , qui est une algèbre graduée à degrés positifs, réduite
\margin{\unsure{\S}}
à $\underline{F}_2$ en degré\deleted{s} 0. C'est d'ailleurs une algèbre de polynômes \inserted{à une infinité de générateurs
$x_i$ ($i$ entier $\unsure{\geq|>}1$\unsure{,} $\neq$ de la forme $2^\nu-1$)} connue ...
\inserted{d'ann\unsure{\deleted{a}e}\unsure{\deleted{d}a}\unsure{\deleted{x}u}xp\unsure{\deleted{d}g}\unsure{\deleted{i}r}\unsure{\deleted{e}a}\unsure{\deleted{e}d}ués}
On a un homomorphisme évident $B.(X)\to H^{\unsure{\deleted{P}L|\deleted{L}P}}(X)$, nul en degrés $>$0, qui
à tout $Z$ sur $X$ associe l'image de $1\in H^o(Z)$ par l'homomorphisme de Gysin
associé à $Z\to X$. Cet homomorphisme par Thom est surjectif (toute clas-%
\margin{en fait $H^\bullet(X)\simeq B.(X)/\Lambda^+B.(X)$}
se de cohomologie provient d'une image de variété compacte ...); \deleted{noter}
On ppouve même qu'il existe un relèvement fonctoriel en $X$
\margin{?}
$H^\bullet(X)\to B.(X)$, (\inserted{il est douteux qu'il soit} compatible avec multiplication) \deleted{tel}tel que l'on ait:

\deleted{a)}\unsure{1)} L'homomorphisme correspondant $H^\bullet(X)\otimes_{\underline{F}_2}\Lambda\simeq B.(X)$ est un iso-%
\margin{?}
morphisme (ce sera vrai pour tout relèvement individuellement ...), \deleted{et}

\deleted{b)}2) [\inserted{Quillen ignore si} \deleted{Le}le relèvement est compatible aussi avec les hom. de Gysin.]

Ainsi, la catégorie $\underline{B}$ admet un foncteur pleinement fidèle dans la
catégorie des modules gradués libres de type fini sur $\Lambda$, les modules
de type fini qu'on obtient ainsi ayant des générateurs à degrés $<$ 0 et
satisfaisant (pour $X$ connexe i.e. quand il y a un seul générateur de
degré zéro) la dualité de Poincaré. Quand on prend la catég\unsure{\deleted{b}o|\deleted{o}b}rie quasi-%
abélienne-\inserted{$\unsure{\widetilde{B}|B}$}associée à $\underline{B}$, on trouve donc tous les modules gradués libres
de type fini à générateurs \deleted{négatifs} en degrés $\leq$ 0 seulement. C'est une

\margin{45}
